\subsection{Parameter Sensitivity and Computational Analysis}
\label{subsec:parameter_sensitivity}

This subsection investigates the effects of the neighborhood sample count
$n$ and the sample spacing $h$ on the three proposed gradient estimators.
The practical implementation uses a separable one-dimensional Taylor model
independently along the horizontal and vertical axes. Therefore, $n$ denotes
the number of symmetric samples per axis, while $h$ determines the distance
between consecutive sampling locations. The polynomial order is fixed to
$p=3$ throughout this analysis.

\subsubsection{Effect of the Neighborhood Sample Count}

To study the influence of $n$, the spacing parameter is fixed to $h=1$, and
$n\in\{4,8,16,32\}$ is evaluated. A representative qualitative comparison is
shown in Fig.~\ref{fig:n_visual_comparison}, while the corresponding numerical
results are summarized in Table~\ref{tab:parameter_summary}.

Increasing $n$ generally improves the noise stability of the estimated
gradient magnitude. In particular, the gradient RMSE decreases substantially
as the number of samples increases. The correntropy-based method also exhibits
a consistent reduction in HOG relative error and an increase in cosine
similarity for larger values of $n$. However, this improvement is accompanied
by a considerable increase in computational cost, especially for the
correntropy-based estimator.

For the squared-loss methods, a lower gradient RMSE does not always lead to a
more stable HOG descriptor. This observation indicates that increasing the
support of the polynomial fit may improve gradient consistency while modifying
local orientation information. Consequently, the gradient and descriptor
metrics are reported separately.

\begin{figure}[t]
	\centering
	\includegraphics[width=\linewidth]{figures/n_visual_comparison}
	\caption{Representative HOG results for different values of the symmetric
		sample count $n$, with $h=1$ and polynomial order $p=3$.}
	\label{fig:n_visual_comparison}
\end{figure}

\subsubsection{Effect of the Sample Spacing}

The influence of $h$ is evaluated by fixing $n=8$ and considering
$h\in\{1,2,3,4\}$. Figure~\ref{fig:h_visual_comparison} presents a representative
visual comparison.

Within the tested interval, increasing $h$ improves the matched-clean
stability metrics under salt-and-pepper noise. The HOG cosine similarity
increases, whereas the HOG relative error and gradient RMSE decrease. The
correntropy-based method consistently provides the most stable descriptor.
The value $h=4$ achieves the best measured noise stability, while $h=3$ is a
more conservative choice when edge localization and numerical conditioning
are also considered. Since $h=4$ is the boundary of the tested range, it is
not claimed to be a global optimum.

\begin{figure}[t]
	\centering
	\includegraphics[width=\linewidth]{figures/h_visual_comparison}
	\caption{Representative HOG results for different values of the sample
		spacing $h$, with $n=8$ and polynomial order $p=3$.}
	\label{fig:h_visual_comparison}
\end{figure}

\subsubsection{Quantitative Comparison}

Table~\ref{tab:parameter_summary} summarizes representative configurations
using gradient RMSE, HOG relative error, cosine similarity, and execution time.
The reported errors compare each noisy output with the clean output generated
using the same estimator parameters. Therefore, these metrics quantify noise
stability rather than absolute derivative accuracy.

\begin{table}[t]
	\centering
	\caption{Representative parameter configurations and their accuracy--runtime
		trade-offs. The numerical entries should be populated automatically from the
		experiment CSV file.}
	\label{tab:parameter_summary}
	\resizebox{\linewidth}{!}{
		\begin{tabular}{lccccccc}
			\hline
			Method & $n$ & $h$ & Noise & Gradient RMSE &
			HOG error & Cosine & Time (s) \\
			\hline
			Direct squared loss    & 4  & 3 & SP 5\%  & -- & -- & -- & -- \\
			Direct squared loss    & 8  & 3 & SP 10\% & -- & -- & -- & -- \\
			Correntropy loss       & 16 & 1 & SP 5\%  & -- & -- & -- & -- \\
			Correntropy loss       & 32 & 1 & SP 10\% & -- & -- & -- & -- \\
			Correntropy loss       & 8  & 4 & SP 5\%  & -- & -- & -- & -- \\
			\hline
	\end{tabular}}
\end{table}

\subsubsection{Computational Complexity}

Figure~\ref{fig:runtime_analysis} shows the execution time as a function of
the image size and the sample count. The non-iterative squared-loss estimator
is consistently the fastest method. The iterative squared-loss estimator
requires additional computation while producing results close to the direct
solution. The correntropy-based estimator achieves the strongest robustness
against impulsive noise, but its repeated weighted least-squares updates
introduce a substantial runtime penalty.

The computational cost increases with both the image dimensions and the
number of neighborhood samples. Therefore, large values of $n$ should only be
used when robustness has priority over inference time. A logarithmic time axis
is used because the execution times of the compared methods differ by more
than one order of magnitude.

\begin{figure}[t]
	\centering
	\includegraphics[width=\linewidth]{figures/runtime_analysis}
	\caption{Execution-time comparison for the three proposed methods over
		different image sizes and neighborhood sample counts.}
	\label{fig:runtime_analysis}
\end{figure}

\subsubsection{Recommended Configurations}

The results reveal that no single configuration is optimal for all practical
requirements. Three operating profiles are therefore recommended:

\begin{itemize}
	\item \textbf{Efficiency-oriented configuration:}
	direct squared loss with a small neighborhood, such as $n=4$ or $n=8$.
	This configuration provides the lowest execution time.
	
	\item \textbf{Balanced configuration:}
	correntropy loss with approximately $n=16$ and a moderate spacing such as
	$h=3$. This setting provides a practical compromise between descriptor
	stability and computational cost.
	
	\item \textbf{Robustness-oriented configuration:}
	correntropy loss with $n=32$ and $h=1$, which achieves the strongest
	measured robustness among the jointly observed configurations.
\end{itemize}

In the independent spacing sweep with $n=8$, $h=4$ provides the best tested
noise stability. However, the present one-factor-at-a-time experiments do not
establish $(n,h)=(32,4)$ as a joint optimum. Overall, Method~I is preferred
when computational efficiency is dominant, Method~II provides only limited
improvement over Method~I for additional cost, and Method~III is preferred
when robustness to impulsive noise is the primary objective.